\section{Strategi Membangun Bukti Formal Berantai}

Kompetensi inferensi yang sejati teruji ketika kita dihadapkan pada puluhan premis dan diminta membangun bukti berantai menuju kesimpulan \cite{rosen2019}.

Seni merekatkan puluhan bongkah bata deduksi inilah yang dipanggil dengan \textbf{Seni Bukti Formal Berantai} (\textit{Formal Proof of Validity}).

\subsection{Prosedur Standar Menyusun Bukti}

Setiap langkah dalam bukti berantai harus disertai justifikasi: baris mana yang dirujuk dan aturan inferensi mana yang diterapkan.

\begin{contoh}[Studi Kasus Argumen Majemuk]
Diberikan tiga premis berikut:
\begin{enumerate}
    \item \textit{"Apabila mahasiswa ini tidak getol merangkum diktat logika ($\neg p$), niscaya kepalanya rontok digebuk Ujian Kuis besok ($q$)"}
    \item \textit{"Asal usut punya sadar, Kepala mahasiswa itu ternyata sakti kebal, BUKAN digebuk malang oleh kuis esok hari ($\neg q$)"}
    \item \textit{"Seandainya narasi menyatakan si malang lolos dari palu kuis esok ($\neg q$), Maka ini mengisyaratkan si cemerlang ini niscaya berpredikat mahasiswa kelewat Cerdas ($r$)"}
\end{enumerate}
\textbf{Tujuan:} Buktikan ($r \land p$) dari ketiga premis di atas.
\end{contoh}

\subsection{Alur Langkah Pengerjaan Pembuktian Bertingkat}

\textbf{Langkah 0: Formulasi Simbolik} \\
Terjemahkan ke bentuk formal: 
\begin{itemize}
    \item 1. $\neg p \rightarrow q$ \quad (Premis 1 asal)
    \item 2. $\neg q$ \quad (Premis 2 saksi)
    \item 3. $\neg q \rightarrow r$ \quad (Premis 3 fakta)
    \item \textbf{Konklusi yang dibuktikan}: \quad $r \land p$
\end{itemize}

\textbf{Langkah Ekskusi Ekstraksi Bertahap:}

Tabel bukti formal:
\begin{table}[htbp]
\centering
\small
\begin{tabular}{|p{0.9cm}|p{3.2cm}|p{6.8cm}|}
\hline
\textbf{Baris} & \textbf{Ekspresi} & \textbf{Justifikasi/Alasan} \\
\hline
1. & $\neg p \rightarrow q$ & Premis 1 \\
2. & $\neg q$                 & Premis 2 \\
3. & $\neg q \rightarrow r$   & Premis 3 \\
\hline
4. & $\neg (\neg p)$          & Modus Tollens (1), (2) \\
5. & $p$                      & Double Negation (4) \\
6. & $r$                      & Modus Ponens (3), (2) \\
\hline
7. & \textbf{$\therefore r \land p$} & Konjungsi (5), (6) \\
\hline
\end{tabular}
\end{table}

\textbf{Kesimpulan}: Konklusi ($r \land p$) terbukti valid dari ketiga premis. Argumen valid.

Metode pembuktian baris-demi-baris dengan justifikasi ini disebut \textbf{Deduksi Formal}. Metode ini digunakan dalam verifikasi perangkat lunak dan model checking \cite{wiki_first_order_logic}.
